\documentclass[a4paper,12pt]{article}
\usepackage[T1]{fontenc}
\usepackage[utf8]{inputenc}
\usepackage{xcolor}
\usepackage{geometry}
\geometry{hmargin=0.5pt ,vmargin=0.5pt}
\usepackage{babel}
\usepackage{graphicx}
\usepackage{graphics}
\title{\textbf{\underline{RECUEIL D'EXERCICE CORRIGE THERMODYNAMIQUE}}} 
\author{MIWADINOU Clement  \\ Maitre Conference dans les universites de CAMES }


\begin{document} 	 
	

	\maketitle 
	
\underline{ENS NATITINGOU } \hfill          \textbf{PC 2 - MI 2} 

\vspace{0.5cm}
	\rule{\linewidth}{0.5pt}  

	\vspace{0.5cm}  
	\begin{center}
			\includegraphics[width=2cm, height=2cm]{"C:/Users/HP/OneDrive/Desktop/MI 2 de natitingou/1762539740531"}%je dois bucher
		\caption{}
		\label{fig:1762539740531}
	\end{center} 
	\vspace{0.5cm} 
	\textbf{PARTIE I } \\
	
	ÉNONCÉS DES EXERCICES \\
	\vspace{0.5cm} 
	\textbf{PARTIE II} \\ 
	
	CORRIGES DES EXERCICES  \\ 
	
	\footnote{@ensnati} 
	\newpage
	\vspace{3cm}
\begin{center}
	\Huge PARTIE I	
\end{center} 
\vspace{10cm} 
\begin{center}
		\Huge LES ÉNONCES DES EXERCICES 
\end{center}
	
		
\footnote{thermodynamique}	 

\newpage

	
	\hrule
	\vspace{0.4cm}
	
	\section*{Questions de cours}
	
	1- Définir un gaz parfait et établir à partir d'un raisonnement scientifique cohérent l'équation d'état d'un gaz parfait.
	
	2- Montrer que, dans un diagramme de Clapeyron donnant la pression $P$ en fonction de $V$, la pente d'une adiabatique est $\gamma$ fois plus élevée que celle d'une isotherme. On donnera la signification de $\gamma$.
	
	3- Montrer que pour les diagrammes entropiques la croissance des courbes représentant les isochores est plus forte que celle des isobares.
	
	4- Un moteur ditherme au cours de son fonctionnement échange de la chaleur avec une source froide de température $T_f$ et une source chaude de température $T_c$.
	\begin{itemize}
		\item[a-] Définir une source thermique et décrire le cycle de Carnot.
		\item[b-] Dans un raisonnement scientifique bien cohérent établir l'expression du rendement maximal $\eta$ d'un moteur ditherme.
		\item[c-] Etablir l'expression du rendement $\eta_c$ du cycle de Carnot.
		\item[d-] Comparer $\eta$ et $\eta_c$. Conclure.
	\end{itemize}
	
	\section*{Exercice 1}
	Du dioxygène contenu dans une bouteille de volume $V_1 = 1\,\text{L}$ à $T_1 = 293\,\text{K}$ sous $P_1 = 3\,\text{bars}$ et du dioxyde de carbone contenu dans une autre bouteille de volume $V_{s2} = 3\,\text{L}$ à $T_2 = 323\,\text{K}$ sous $P_2 = 2\,\text{bars}$, sont introduits en totalité dans un réservoir de volume $V = 5\,\text{L}$ initialement vide. La température après mélange est $T = 313\,\text{K}$. Calculer la pression finale $P$ dans le réservoir et les pressions partielles $p_1$ et $p_2$ du dioxygène et de dioxyde de carbone.
	
	\section*{Exercice 2}
	Un thermomètre à mercure, gradué linéairement, est plongé dans la glace fondante ; le mercure affleure à la division $-2$. Dans la vapeur d'eau bouillante, sous $76\,\text{cm}$ de mercure, il affleure à la division $+103$.
	
	1- Dans un bain tiède, le mercure affleure à la division $n = +70$. Déterminer la température $\theta$ du bain, indiquée par ce thermomètre.
	
	2- Plus généralement, déterminer la correction à apporter à la lecture de la division $n$, sous la forme $\theta - n = f(n)$. En déduire la température $\theta$ pour laquelle aucune correction n'est nécessaire.
	
	\section*{Exercice 3}
	Une mole de gaz carbonique obéit à l'équation de Van der Waals :
	\begin{equation}
		\left(p + \frac{a}{v^2}\right)(v-b) = RT
	\end{equation}
	
	1- Exprimer, en fonction des variables indépendantes : volume $v$ et température absolue $T$, les coefficients de dilatation à pression constante $\alpha$ et à volume constant $\beta$.
	
	2- Trouver la relation générale entre le coefficient $\chi$ de compressibilité isotherme, les coefficients $\alpha$ et $\beta$ et la pression $p$ du gaz. En déduire le coefficient $\chi$ du gaz de Van der Waals.
	
	3- Dans le cas où l'on peut négliger la pression interne du gaz, montrer que $\chi = \frac{v-b}{p \cdot v}$.
	
	\section*{Exercice 4}
	1- Les états d'équilibre d'un gaz sont décrits par une équation d'état du type $f(p,v,T) = 0$. Dans la représentation d'Amagat $(pv, p)$, calculer d'une façon générale, la pente $m$ d'une isotherme, en un point donné, en fonction de $\chi$ et $\chi_0$ ($\chi$ et $\chi_0$ sont respectivement les coefficients de compressibilité isotherme du gaz étudié et d'un gaz parfait pris à la même pression). Que peut-on conclure pour les régions où $m$ est positif et pour celles où $m$ est négatif ?
	
	2- Soit l'équation d'état de Van der Waals écrite pour 1 gramme de gaz :
	\begin{equation}
		\left(p + \frac{a}{v^2}\right)(v-b) = rT
	\end{equation}
	Retrouver les valeurs des coefficients $a$, $b$, $r$ en fonction des constantes critiques $p_c$, $v_c$, $T_c$ du gaz.
	
	3- 
	\begin{itemize}
		\item[a)] On pose $pv = y$. Écrire l'équation $f(y,p) = \text{cte}$ d'une isotherme, dans la représentation d'Amagat, pour le gaz considéré.
		\item[b)] Quel est le lieu $(C)$ des points où le gaz suit locally la loi de Mariotte ? Tracer cette courbe.
		\item[c)] Quelles sont les régions où $m$ est respectivement positif et négatif ? Dessiner quelques isothermes dans cette représentation ($T > T_c$).
		\item[d)] Calculer la température $T_1$ à laquelle le gaz suit la loi de Mariotte pour une pression nulle. Comparer $T_1$ à $T_c$.
	\end{itemize}
	Application numérique : Cas de l'azote : $T_c = 126\,\text{K}$.
	
	\section*{Exercice 5}
	Un gaz monoatomique est maintenu à une pression de $1{,}2\,\text{bar}$ et à une température $300\,\text{K}$, dans une enceinte cylindrique de volume $V_i = 1\,\text{L}$, grâce à une masse $M$ posée sur un piston de masse $m_p = 1\,\text{kg}$. Le piston est à une hauteur $h_i = 50\,\text{cm}$. On enlève la masse $M$, ce qui permet au gaz de se détendre, de façon adiabatique, jusqu'à la pression finale d'équilibre $p_f$ ; le volume est alors $V_f$. On désigne par $p_0 = 1\,\text{bar}$ la pression atmosphérique.
	
	1- Calculer la valeur de la masse posée sur le piston et la pression finale $p_f$.
	
	2- Effectuer le bilan énergétique de la transformation. Trouver les rapports $V_f/V_i$ et $T_f/T_i$. \\
	Application numérique : On donne $C_{vm} = 3R/2$ et on admits que l'équation d'état du gaz est celle d'un gaz parfait. Calculer le travail reçu par le gaz.
	
	\section*{Exercice 6}
	Quelle quantité de chaleur faut-il fournir, à pression atmosphérique, à $100\,\text{g}$ de glace, initialement à $268\,\text{K}$ pour la transformer en vapeur d'eau à $383\,\text{K}$ ?
	
	On donne les capacités thermiques massiques de la glace $C_1 = 2{,}09\,\text{J}\cdot\text{g}^{-1}\cdot\text{K}^{-1}$, de l'eau liquide $C_2 = 4{,}18\,\text{J}\cdot\text{g}^{-1}\cdot\text{K}^{-1}$, de la vapeur $C_3 = 1{,}87\,\text{J}\cdot\text{g}^{-1}\cdot\text{K}^{-1}$ et les enthalpies massiques de fusion de la glace $L_f = 335\,\text{kJ}\cdot\text{kg}^{-1}$ et de vaporisation de la vapeur d'eau $L_v = 2257\,\text{kJ}\cdot\text{kg}^{-1}$.
	
	\section*{Exercice 7}
	Un gaz réel suit l'équation de Van der Waals :
	\begin{equation}
		\left(P + \frac{a}{V^2}\right)(V-b) = RT
	\end{equation}
	Il subit la succession de transformations suivantes qui seront supposées quasi-statiques :
	\begin{itemize}
		\item[\textbf{A $\rightarrow$ B :}] dilatation isobare à la pression $P_0$ augmentant le volume de $V_0 = 1\,\text{L}$ à $V_1 = 5\,\text{L}$.
		\item[\textbf{B $\rightarrow$ C :}] compression isotherme à la température de $1000\,\text{K}$ ramenant le volume à $V_0$.
		\item[\textbf{C $\rightarrow$ A :}] détente isochore ramenant le système à son état initial.
	\end{itemize}
	1- Tracer ce cycle dans le diagramme de Watt.\\
	2- Dans les unités du système international, $a = 0{,}14$ et $b = 3{,}22\cdot10^{-5}$. Préciser ces unités.\\
	3- Préciser les valeurs de la température, de la pression et du volume aux sommets du cycle.\\
	4- Calculer les travaux des forces de pression au cours de chaque phase du cycle.\\
	5- Au cours du cycle, le système reçoit-il ou cède-t-il un transfert thermique ? Justifier numériquement.\\
	6- Déterminer les transferts thermiques au cours de chaque phase du cycle.\\
	7- Quelle vérification peut-on faire ? On donne l'expression de son énergie interne :
	\begin{equation}
		U = \frac{3}{2}RT - \frac{a}{V}
	\end{equation}
	
	\section*{Exercice 8}
	A l'intérieur d'une concavité vide dont les parois sont à l'équilibre thermique, il existe des ondes électromagnétiques dont l'intensité et la répartition des fréquences dépend de la température. A une onde électromagnétique on associe des particules "grains d'énergie" appelées photons et on considère que ceux-ci se comportent comme un gaz. Des considérations théoriques amènent à poser l'expression de l'entropie, fonction du volume et de l'énergie :
	\begin{equation}
		S(V,U) = \frac{4}{3}(\sigma_0 V U^3)^{\frac{1}{4}}
	\end{equation}
	où $\sigma_0$ est une constante universelle ; $\sigma_0 = 7{,}56\cdot10^{-16}\,\text{J}\cdot\text{m}^{-3}\cdot\text{K}^{-4}$.
	
	1- En déduire la relation entre l'énergie, le volume et la température, puis la relation entre pression et température.
	
	2- La cavité, de volume $V = 1\,\text{L}$ contient 1 mole de dihydrogène. Pour quelle température la pression de radiation est-elle égale à la pression exercée par les particules matérielles (on supposera que, pour les températures très élevées, les molécules de $H_2$ sont décomposées en protons et en électrons qui se comportent comme des gaz parfaits monoatomiques).
	
	3- Déterminer la capacité thermique à volume constant associée au rayonnement.
	
	\section*{Exercice 9}
	NB: Les questions 1, 2 et 3 de cet exercice sont indépendantes.
	
	1) Un local, de capacité thermique $C = 4\cdot10^3\,\text{kJ}\cdot\text{K}^{-1}$, est initialement à la température de l'air extérieur $T_0 = 305\,\text{K}$. Un climatiseur, fonctionnant de façon cyclique réversible ditherme entre l'air extérieur comme source chaude et le local comme source froide, ramène la température du local à $T_1 = 293\,\text{K}$ en 1 h. Quelle est la puissance électrique moyenne $P$ reçue par le climatiseur ?
	
	2) Dans une centrale thermique, de l'air, supposé parfait, décrit de façon réversible un cycle de moteur de Joule dans le sens 1-2-3-4-1. Il entre dans la turbine, à la pression de $7\,\text{bar}$ et à la température de $1227\,\text{K}$, et sort à la pression de $1\,\text{bar}$. Un compresseur comprime l'air qu'il reçoit, à la température $288\,\text{K}$, de 1 à $7\,\text{bar}$.
	\begin{itemize}
		\item[a)] Représenter le cycle dans le diagramme de Clapeyron.
		\item[b)] Calculer l'efficacité de la machine.
	\end{itemize}
	
	3) Un réfrigérateur fonctionne entre deux sources $\Sigma_c$ et $\Sigma_f$ de températures constantes $T_c = 300\,\text{K}$ et $T_f = 263\,\text{K}$ respectivement. On désigne par $W$ le travail reçu par la machine, $Q_f$ la chaleur qu'elle reçoit de $\Sigma_f$, $Q_c$ la quantité analogue provenant de $\Sigma_c$, et $S_p$ l'entropie produite, pour une masse déterminée de fluide, après une évolution cyclique.
	\begin{itemize}
		\item[a)] On constate que le rapport des quantités de chaleur est relié au rapport des températures par l'équation: $Q_c/Q_f = -k T_c/T_f$, $k$ étant une constante. Montrer que $k$ a un signe déterminé. Trouver l'efficacité $\eta_r$ en fonction des températures et de $k$. Application numérique pour : $k=1,2$.
		\item[b)] Exprimer $\eta_r$ en fonction de $T_f, T_c, Q_f$ et $S_p$. En déduire sa valeur maximale $(\eta_r)_{max}$ pour laquelle $S_p = 0$ et la comparer à $\eta_r$. Conclure.
	\end{itemize}
	
	\section*{Exercice 10}
	Un moteur thermique, utilisant un fluide parfait, décrit un cycle réversible Diesel $A_1A_2A_3A_4A_1$ composé d'une isobare et d'une isochore reliées par deux adiabatiques :
	\begin{itemize}
		\item[$A_1A_2$ :] l'air admis subit une compression adiabatique de l'état initial $A_1(p_1, v_1, T_1)$ à l'état $A_2(p_2, v_2, T_2)$
		\item[$A_2A_3$ :] combustion isobare par injection progressive du carburant entre l'état $A_2$ et l'état $A_3(v_3, T_3)$
		\item[$A_3A_4$ :] l'injection cesse en $A_3$ et le mélange subit une détente adiabatique jusqu'à l'état $A_4(v_4=v_1, T_4)$.
		\item[$A_4A_1$ :] refroidissement isochore.
	\end{itemize}
	
	1- Représenter le cycle Diesel sur un diagramme $(p, v)$ et sur un diagramme $(T, S)$ entropique.\\
	2- Exprimer le rendement du cycle Diesel, en fonction :
	\begin{itemize}
		\item[a)] des températures $T_1$, $T_2$, $T_3$ et $T_4$ et du rapport $\gamma$ des chaleurs massiques du mélange gazeux.
		\item[b)] du "taux de compression" $x = \frac{v_1}{v_2}$, du "taux de détente" $y = \frac{v_1}{v_3}$ (ou $y = \frac{v_4}{v_3}$) et du rapport $\gamma$.
	\end{itemize}
	3- Une automobile à moteur Diesel possède les caractéristiques suivantes : taux de compression $x=21$ ; taux de détente $y=7$. A la vitesse maximale du véhicule : $v = 147\,\text{km/h}$ correspondant à $N = 450\,\text{tours/minute}$, la consommation est $c = 8\,\text{litres}$ de carburant (gas-oil) aux $100\,\text{km}$. Le gaz-oil a une masse volumique $\rho = 0{,}8\,\text{kg/l}$ et un pouvoir calorifique $q = 46{,}8\,\text{kJ/g}$. Déterminer (on donne $\gamma = 1{,}4$) :
	\begin{itemize}
		\item[a)] le rendement théorique de ce moteur Diesel.
		\item[b)] la masse de carburant injectée à chaque cycle, à vitesse maximale.
		\item[c)] la puissance maximale de ce moteur Diesel, supposé idéal.
	\end{itemize}
	
	\section*{Exercice 11}
	Dans un moteur à explosion, un fluide supposé parfait ($n$ moles) décrit le cycle Beau de Rochas, appelé aussi Cycle d'Otto, composé de deux adiabatiques et deux isochores :
	\begin{itemize}
		\item - compression adiabatique de l'état $A_1(p_1, v_1, T_1)$ à l'état $A_2(p_2, v_2, T_2)$
		\item - échauffement isochore de l'état $A_2(p_2, T_2)$ à l'état $A_3(p_3, T_3)$
		\item - détente adiabatique de l'état $A_3(p_3, T_3)$ à l'état $A_4(p_4, T_4)$
		\item - refroidissement isochore qui ramène le fluide à l'état initial.
	\end{itemize}
	1- Représenter le cycle d'Otto dans le diagramme $(p, v)$ et dans le diagramme entropique $(T, S)$.\\
	2- Exprimer le rendement théorique $r$ de ce cycle :
	\begin{itemize}
		\item[a)] en fonction des températures $T_1$, $T_2$, $T_3$ et $T_4$ ;
		\item[b)] puis en fonction du rapport volumique $x = \frac{v_1}{v_2}$ et du rapport $\gamma = \frac{C_p}{C_v}$ des chaleurs massiques du fluide.
	\end{itemize}
	3- Le piston du cylindre où évolue l'air ($\gamma = 1{,}4$) a une course utile $l = 10\,\text{cm}$, une section $S = 50\,\text{cm}^2$ et emprisonne un volume d'air de $100\,\text{cm}^3$ en fin de compression. Calculer :
	\begin{itemize}
		\item[a)] le rendement théorique du cycle d'Otto.
		\item[b)] le travail fourni au cours d'un cycle, si l'air est admis sous l'atmosphère à $300\,\text{K}$, et si la température maximale atteinte est $900\,\text{K}$.
	\end{itemize}
	4- Pour quelle valeur du rapport volumétrique, le moteur fonctionnant suivant le cycle d'Otto entre les températures $300\,\text{K}$ et $900\,\text{K}$ a-t-il le même rendement qu'une machine réversible fonctionnant suivant le cycle de Carnot entre les mêmes températures ?   
	\\ 
	
	\footnote{thermodynamique} 
	
	\newpage 
	
 \centering	\Huge PARTIE 2  
 
 \vspace{5cm}
 
 \centering{ \Huge CORRIGES DES EXERCICES}   \\ 
 
 \newpage 
   
 
 	\hrule
 	\vspace{0.4cm}
 	
 
 		\large \textit{Recueil des Solutions des 11 Exercices}
 	
 	\vspace{0.5cm}
 	
 	% ==========================================
 	\section*{Exercice 1 }
 	% ==========================================
 	\subsection*{1. Calcul de la pression finale $P$ dans le réservoir}
 	D'après la loi des gaz parfaits, pour le dioxygène ($O_2$, constituant 1) et le dioxyde de carbone ($CO_2$, constituant 2) :
 	\[ n_1 = \frac{P_1 V_{s1}}{R T_1} \quad \text{et} \quad n_2 = \frac{P_2 V_{s2}}{R T_2} \]
 	Après mélange dans le volume total $V = 5\text{ L}$ à la température $T = 313\text{ K}$ ($40^\circ\text{C}$), la pression finale $P$ est :
 	\[ P = \frac{(n_1 + n_2)RT}{V} = \frac{T}{V} \left( \frac{P_1 V_{s1}}{T_1} + \frac{P_2 V_{s2}}{T_2} \right) \]
 	En convertissant les volumes en $\text{m}^3$ ($1\text{ L} = 10^{-3}\text{ m}^3$) :
 	\[ P = \frac{313}{5 \cdot 10^{-3}} \left( \frac{3 \cdot 10^5 \times 1 \cdot 10^{-3}}{293} + \frac{2 \cdot 10^5 \times 3 \cdot 10^{-3}}{323} \right) \approx 1,80 \cdot 10^5\text{ Pa} \quad (1,80\text{ bar}) \]
 	
 	\subsection*{2. Pressions partielles $p_1$ et $p_2$}
 	Les pressions partielles dans le mélange final à $T = 313\text{ K}$ s'expriment par :
 	\[ p_1 = \frac{n_1 RT}{V} = \frac{P_1 V_{s1}}{T_1} \frac{T}{V} = \frac{3 \cdot 10^5 \times 10^{-3}}{293} \times \frac{313}{5 \cdot 10^{-3}} \approx 6,41 \cdot 10^4\text{ Pa} \quad (0,64\text{ bar}) \]
 	\[ p_2 = \frac{n_2 RT}{V} = \frac{P_2 V_{s2}}{T_2} \frac{T}{V} = \frac{2 \cdot 10^5 \times 3 \cdot 10^{-3}}{323} \times \frac{313}{5 \cdot 10^{-3}} \approx 1,16 \cdot 10^5\text{ Pa} \quad (1,16\text{ bar}) \]
 	Vérification de la loi de Dalton : $p_1 + p_2 = 0,641 + 1,163 = 1,804\text{ bar} = P$.
 	
 	\newpage
 	
 	% ==========================================
 	\section*{Exercice 2 :}
 	% ==========================================
 	La température empirique $\theta$ mesurée par la hauteur de la colonne de mercure est modélisée par une relation linéaire avec la division $n$ :
 	\[ \theta = a \cdot n + b \]
 	
 	\subsection*{1. Détermination des constantes et de la température du bain}
 	En utilisant les points fixes d'étalonnage de l'eau :
 	\begin{itemize}
 		\item Glace fondante ($\theta_1 = 0^\circ\text{C}$) : $n_1 = -2 \implies -2a + b = 0 \implies b = 2a$
 		\item Vapeur d'eau bouillante ($\theta_2 = 100^\circ\text{C}$) : $n_2 = 103 \implies 103a + b = 100$
 	\end{itemize}
 	En substituant $b$ :
 	\[ 105a = 100 \implies a = \frac{20}{21} \approx 0,952 \quad \text{et} \quad b = \frac{40}{21} \approx 1,905 \]
 	La loi d'étalonnage est donc : $\theta = \frac{20}{21}n + \frac{40}{21}$.
 	Pour une lecture de division $n = 70$ :
 	\[ \theta = \frac{20}{21}(70) + \frac{40}{21} = \frac{1440}{21} \approx 68,57^\circ\text{C} \]
 	
 	\subsection*{2. Température sans correction}
 	Aucune correction n'est nécessaire si la valeur lue est égale à la vraie valeur ($\theta = n$) :
 	\[ n = \frac{20}{21}n + \frac{40}{21} \implies n \left(1 - \frac{20}{21}\right) = \frac{40}{21} \implies \frac{n}{21} = \frac{40}{21} \implies n = 40 \]
 	La température correspondante est donc $\theta = 40^\circ\text{C}$.
 	
 	% ==========================================
 	\section*{Exercice 3 : }
 	% ==========================================
 	L'équation d'état pour $1\text{ mole}$ est :
 	\[ \left(P + \frac{a}{v^2}\right)(v-b) = RT \implies P(v,T) = \frac{RT}{v-b} - \frac{a}{v^2} \]
 	
 	\subsection*{1. Expressions des coefficients thermoélastiques}
 	Par définition, les coefficients $\alpha$, $\beta$ et la compressibilité isotherme $\chi_T$ sont :
 	\[ \alpha = \frac{1}{v} \left(\frac{\partial v}{\partial T}\right)_P, \quad \beta = \frac{1}{P} \left(\frac{\partial P}{\partial T}\right)_v, \quad \chi_T = -\frac{1}{v} \left(\frac{\partial v}{\partial P}\right)_T \]
 	Calcul de $\beta$ :
 	\[ \left(\frac{\partial P}{\partial T}\right)_v = \frac{R}{v-b} \implies \beta = \frac{R}{P(v-b)} \]
 	Pour obtenir $\alpha$ et $\chi_T$, on utilise la relation cyclique $\left(\frac{\partial v}{\partial T}\right)_P \left(\frac{\partial T}{\partial P}\right)_v \left(\frac{\partial P}{\partial v}\right)_T = -1$ :
 	\[ \left(\frac{\partial P}{\partial v}\right)_T = -\frac{RT}{(v-b)^2} + \frac{2a}{v^3} \]
 	On en déduit :
 	\[ \alpha = \frac{1}{v} \frac{\left(\frac{\partial P}{\partial T}\right)_v}{-\left(\frac{\partial P}{\partial v}\right)_T} = \frac{R v^2 (v-b)}{R T v^3 - 2a(v-b)^2} \]
 	\[ \chi_T = -\frac{1}{v \left(\frac{\partial P}{\partial v}\right)_T} = \frac{v^2 (v-b)^2}{R T v^3 - 2a(v-b)^2} \]
 	
 	\subsection*{2. Relation générale entre $\chi_T$, $\alpha$, $\beta$ et la pression $P$}
 	On constate en faisant le rapport des définitions que :
 	\[ \frac{\alpha}{\chi_T} = \left(\frac{\partial P}{\partial T}\right)_v = P\beta \implies \chi_T = \frac{\alpha}{P\beta} \]
 	
 	\newpage
 	
 	% ==========================================
 	\section*{Exercice 4 : }
 	% ==========================================
 	\subsection*{1. Pente isotherme d'Amagat $m = \left(\frac{\partial(Pv)}{\partial P}\right)_T$}
 	Le produit $Pv$ pour un gaz réel s'exprime à l'aide des coefficients de compressibilité. On peut exprimer $v$ sous la forme d'un développement en pression $Pv = RT + B(T)P + C(T)P^2 + \dots$ :
 	\[ m = \left(\frac{\partial(Pv)}{\partial P}\right)_T = B(T) \]
 	Si $m > 0$, le fluide est moins compressible que le gaz parfait (les forces de répulsion à courte portée dominent). Si $m < 0$, il est plus compressible (les forces d'attraction à moyenne portée dominent).
 	
 	\subsection*{2. Constantes critiques pour l'équation de Van der Waals}
 	Au point critique $(P_c, v_c, T_c)$, la courbe isotherme présente un point d'inflexion à tangente horizontale :
 	\[ \left(\frac{\partial P}{\partial v}\right)_T = 0 \implies \frac{RT_c}{(v_c-b)^2} = \frac{2a}{v_c^3} \quad (1) \]
 	\[ \left(\frac{\partial^2 P}{\partial v^2}\right)_T = 0 \implies \frac{2RT_c}{(v_c-b)^3} = \frac{6a}{v_c^4} \quad (2) \]
 	En divisant (2) par (1) :
 	\[ \frac{2}{v_c - b} = \frac{3}{v_c} \implies v_c = 3b \]
 	En injectant $v_c$ dans (1), on trouve :
 	\[ T_c = \frac{8a}{27Rb} \quad \text{et} \quad P_c = \frac{a}{27b^2} \]
 	
 	\subsection*{3. Représentation d'Amagat $Pv = f(P)$}
 	\begin{itemize}
 		\item \textbf{a)} Équation : $Pv = RT + P\left(b - \frac{a}{RT}\right)$ (en première approximation pour $v \gg b$).
 		\item \textbf{b)} Le lieu géométrique des points où le gaz suit localement la loi de Mariotte correspond à la courbe d'annulation de la pente : $m = 0 \implies T_B = \frac{a}{Rb}$ (Température de Boyle-Mariotte).
 		\item \textbf{c)} Pour $T < T_B$, la pente est d'abord négative puis positive. Pour $T > T_B$, elle est toujours positive.
 		\item \textbf{d)} Température d'inversion pour l'Azote ($T_c = 126\text{ K}$) : $T_i \approx 3,3 T_c \approx 415\text{ K}$.
 	\end{itemize}
 	
 	% ==========================================
 	\section*{Exercice 5 : }
 	% ==========================================
 	\subsection*{1. Masse $M$ et pression finale $P_f$}
 	À l'équilibre initial (piston de masse $m_p$, section $S$, hauteur $h_i = 50\text{ cm} = 0,5\text{ m}$) :
 	\[ P_i = P_0 + \frac{(M+m_p)g}{S} \]
 	Le volume initial est $V_i = 1\text{ L} = 10^{-3}\text{ m}^3 \implies S = \frac{V_i}{h_i} = 2\cdot 10^{-3}\text{ m}^2$.
 	\[ M = \frac{S(P_i - P_0)}{g} - m_p = \frac{2\cdot 10^{-3}(1,2\cdot 10^5 - 1\cdot 10^5)}{9,81} - 1 \approx 3,08\text{ kg} \]
 	Lorsque l'on retire la masse $M$ brusquement, la pression d'équilibre finale est assurée par le seul piston :
 	\[ P_f = P_0 + \frac{m_p g}{S} = 10^5 + \frac{1 \times 9,81}{2\cdot 10^{-3}} \approx 1,05 \cdot 10^5\text{ Pa} \quad (1,05\text{ bar}) \]
 	
 	\subsection*{2. Bilan énergétique de la transformation irréversible}
 	La paroi étant adiabatique ($Q = 0$). Le premier principe s'écrit :
 	\[ \Delta U = W \implies n C_{vm}(T_f - T_i) = -P_f(V_f - V_i) \]
 	Pour un gaz monoatomique, $C_{vm} = \frac{3}{2}R$. En utilisant $PV = nRT$, on a :
 	\[ \frac{3}{2}(P_f V_f - P_i V_i) = -P_f(V_f - V_i) \implies \frac{5}{2}P_f V_f = \frac{3}{2}P_i V_i + P_f V_i \]
 	\[ \frac{V_f}{V_i} = \frac{3 P_i + 2 P_f}{5 P_f} = \frac{3(1,2) + 2(1,05)}{5(1,05)} = \frac{5,7}{5,25} \approx 1,086 \]
 	\[ \frac{T_f}{T_i} = \frac{P_f V_f}{P_i V_i} = \left(\frac{1,05}{1,2}\right) \times 1,086 \approx 0,950 \]
 	Le travail reçu par le gaz est :
 	\[ W = -P_f(V_f - V_i) = -1,05\cdot 10^5 \times 10^{-3} \times (1,086 - 1) \approx -9,03\text{ J} \]
 	
 	\newpage
 	
 	% ==========================================
 	\section*{Exercice 6 : }
 	% ==========================================
 	Calcul de la quantité de chaleur nécessaire pour transformer $m = 100\text{ g}$ de glace de $T_1 = 268\text{ K}$ ($-5^\circ\text{C}$) à de la vapeur à $T_2 = 383\text{ K}$ ($110^\circ\text{C}$) :
 	\[ Q = Q_1 + Q_2 + Q_3 + Q_4 + Q_5 \]
 	\begin{itemize}
 		\item Chauffage glace de $268\text{ K}$ à $273\text{ K}$ : $Q_1 = m C_1 \Delta T = 0,1 \times 2,09\cdot 10^3 \times 5 = 1045\text{ J}$
 		\item Fusion de la glace à $273\text{ K}$ : $Q_2 = m L_f = 0,1 \times 335\cdot 10^3 = 33500\text{ J}$
 		\item Chauffage de l'eau liquide de $273\text{ K}$ à $373\text{ K}$ : $Q_3 = m C_2 \Delta T = 0,1 \times 4,18\cdot 10^3 \times 100 = 41800\text{ J}$
 		\item Vaporisation de l'eau à $373\text{ K}$ : $Q_4 = m L_v = 0,1 \times 2257\cdot 10^3 = 225700\text{ J}$
 		\item Chauffage de la vapeur de $373\text{ K}$ à $383\text{ K}$ : $Q_5 = m C_3 \Delta T = 0,1 \times 1,87\cdot 10^3 \times 10 = 1870\text{ J}$
 	\end{itemize}
 	\[ Q_{\text{total}} = 1045 + 33500 + 41800 + 225700 + 1870 = 303915\text{ J} \approx 303,9\text{ kJ} \]
 	
 	% ==========================================
 	\section*{Exercice 7 : }
 	% ==========================================
 	Données : $a = 0,14\text{ Pa}\cdot\text{m}^6\cdot\text{mol}^{-2}$, $b = 3,22\cdot 10^{-5}\text{ m}^3\cdot\text{mol}^{-1}$.
 	Le cycle $A \rightarrow B \rightarrow C \rightarrow A$ s'effectue sur $1\text{ mole}$ :
 	\begin{itemize}
 		\item $A \rightarrow B$ : Isobare à $P_0$ de $V_0 = 1\text{ L}$ à $V_1 = 5\text{ L}$.
 		\item $B \rightarrow C$ : Isotherme à $T_{\text{max}} = 1000\text{ K}$ de $V_1$ à $V_0$.
 		\item $C \rightarrow A$ : Isochore ramenant le système à l'état initial.
 	\end{itemize}
 	
 	\subsection*{1. Représentation du cycle (Watt)}
 	Dans un diagramme $(P,V)$, le point $A$ est à basse température, l'isobare $AB$ s'étend vers la droite, l'isotherme $BC$ remonte vers la gauche à petit volume, et l'isochore $CA$ redescend verticalement.
 	
 	\subsection*{2. Calcul du travail pour chaque étape}
 	\begin{itemize}
 		\item $W_{AB} = -P_0(V_1 - V_0) = -P_0(4\cdot 10^{-3})\text{ J}$
 		\item $W_{BC} = -RT_{\text{max}} \ln\left(\frac{V_0-b}{V_1-b}\right) - a\left(\frac{1}{V_0} - \frac{1}{V_1}\right)$
 		\[ W_{BC} = -8,314 \times 1000 \times \ln\left(\frac{10^{-3} - 3,22\cdot 10^{-5}}{5\cdot 10^{-3} - 3,22\cdot 10^{-5}}\right) - 0,14\left(\frac{1}{10^{-3}} - \frac{1}{5\cdot 10^{-3}}\right) \]
 		\[ W_{BC} = -8314 \times (-1,634) - 112 = 13585 - 112 = 13473\text{ J} \]
 		\item $W_{CA} = 0$ (Isochore)
 	\end{itemize}
 	
 	% ==========================================
 	\section*{Exercice 8 : s}
 	% ==========================================
 	L'entropie d'une concavité est donnée par : $S(V,U) = \frac{4}{3}(\sigma_0 V U^3)^{1/4}$.
 	
 	\subsection*{1. Relation entre $U, V, T$ et équation d'état $P(V,T)$}
 	Par l'identité fondamentale $dU = TdS - PdV \implies dS = \frac{1}{T}dU + \frac{P}{T}dV$, on a :
 	\[ \frac{1}{T} = \left(\frac{\partial S}{\partial U}\right)_V = \frac{4}{3} \times \frac{1}{4} (\sigma_0 V)^{1/4} U^{-3/4} = \frac{1}{3} \left(\frac{\sigma_0 V}{U^3}\right)^{1/4} \implies U(V,T) = \sigma_0 T^4 V \]
 	Pour la pression $P$ :
 	\[ \frac{P}{T} = \left(\frac{\partial S}{\partial V}\right)_U = \frac{1}{3} \left(\frac{\sigma_0 U^3}{V^3}\right)^{1/4} \implies P = \frac{T}{3} \left(\frac{\sigma_0 (\sigma_0 T^4 V)^3}{V^3}\right)^{1/4} = \frac{1}{3}\sigma_0 T^4 \]
 	
 	\subsection*{2. Pression de radiation pour 1 mole de $H_2$ décomposé à $V = 1\text{ L}$}
 	À très haute température, $H_2 \rightarrow 2H \rightarrow 2p^+ + 2e^-$, soit $4\text{ moles}$ de gaz parfait par mole initiale de $H_2$.
 	\[ P_{\text{matière}} = \frac{4 R T}{V} \quad \text{et} \quad P_{\text{rad}} = \frac{1}{3}\sigma_0 T^4 \]
 	Égalité : $\frac{1}{3}\sigma_0 T^4 = \frac{4 R T}{V} \implies T = \left(\frac{12 R}{\sigma_0 V}\right)^{1/3}$.
 	\[ T = \left(\frac{12 \times 8,314}{7,56\cdot 10^{-16} \times 10^{-3}}\right)^{1/3} \approx 5,09 \cdot 10^6\text{ K} \]
 	
 	\subsection*{3. Capacité thermique $C_v$}
 	\[ C_v = \left(\frac{\partial U}{\partial T}\right)_V = 4 \sigma_0 T^3 V \]
 	
 	\newpage
 	
 	% ==========================================
 	\section*{Exercice 9 : }
 	% ==========================================
 	\subsection*{1. Climatiseur cyclique ditherme}
 	En $1\text{ h}$, la chaleur extraite de la pièce (source froide à $T_1 = 293\text{ K}$) de capacité thermique $C = 4\cdot 10^3\text{ kJ}\cdot\text{K}^{-1}$ pour la refroidir de $\Delta T = T_0 - T_1 = 12\text{ K}$ est :
 	\[ Q_1 = C (T_1 - T_0) = 4\cdot 10^6 \times (293 - 305) = -4,8\cdot 10^7\text{ J} \]
 	Pour un climatiseur réversible ($S_{\text{créée}} = 0$) :
 	\[ \frac{Q_0}{T_0} + \frac{Q_1}{T_1} = 0 \implies Q_0 = -Q_1 \frac{T_0}{T_1} = 4,8\cdot 10^7 \times \frac{305}{293} \approx 4,996 \cdot 10^7\text{ J} \]
 	Le travail électrique reçu est :
 	\[ W = - (Q_0 + Q_1) = -(4,996\cdot 10^7 - 4,8\cdot 10^7) \approx 1,96\cdot 10^6\text{ J} \]
 	La puissance moyenne nécessaire en $1\text{ h}$ ($3600\text{ s}$) est :
 	\[ P = \frac{W}{\Delta t} = \frac{1,96\cdot 10^6}{3600} \approx 544,4\text{ W} \]
 	
 	\subsection*{2. Turbine à gaz (Cycle de Joule)}
 	a) Le cycle comporte deux isentropiques ($1\rightarrow2$ et $3\rightarrow4$) et deux isobares ($2\rightarrow3$ et $4\rightarrow1$).
 	b) L'efficacité (rendement) d'un cycle de Joule parfait s'exprime par :
 	\[ \eta = 1 - \left(\frac{P_1}{P_2}\right)^{\frac{\gamma-1}{\gamma}} = 1 - \left(\frac{1}{7}\right)^{\frac{0,4}{1,4}} \approx 0,427 \quad (42,7\%) \]
 	
 	\subsection*{3. Réfrigérateur réversible}
 	a) $Q_c/Q_f = -k T_c/T_f \implies k = 1$ pour un cycle parfaitement réversible d'après le théorème de Carnot ($Q_c/T_c + Q_f/T_f = 0$).
 	L'efficacité est : $\eta_r = \frac{Q_f}{W} = \frac{Q_f}{-(Q_c+Q_f)} = \frac{1}{-\frac{Q_c}{Q_f}-1} = \frac{1}{k\frac{T_c}{T_f} - 1}$.
 	Pour $k=1$, $\eta_{r} = \frac{263}{300-263} \approx 7,11$. Pour $k=2$, $\eta_{r} \approx 0,78$.
 	b) Avec l'entropie créée $S_p$ :
 	\[ \frac{Q_c}{T_c} + \frac{Q_f}{T_f} + S_p = 0 \implies Q_c = -T_c \left(\frac{Q_f}{T_f} + S_p\right) \]
 	\[ \eta_r = \frac{Q_f}{-Q_c - Q_f} = \frac{Q_f}{T_c\left(\frac{Q_f}{T_f} + S_p\right) - Q_f} = \frac{1}{\frac{T_c}{T_f} - 1 + \frac{T_c S_p}{Q_f}} \]
 	Si $S_p = 0$, on retrouve l'efficacité maximale de Carnot $(\eta_r)_{\text{max}} = \frac{T_f}{T_c - T_f}$. Comme $S_p > 0$, alors $\eta_r < (\eta_r)_{\text{max}}$.
 	
 	% ==========================================
 	\section*{Exercice 10 : }
 	% ==========================================
 	\subsection*{1. Représentations}
 	Le cycle se trace dans le diagramme $(P,V)$ avec deux adiabatiques encadrant une isobare en haut ($2\rightarrow 3$) et une isochore à droite ($4\rightarrow 1$).
 	
 	\subsection*{2. Rendement thermique}
 	\[ \eta = 1 + \frac{Q_{\text{rejetée}}}{Q_{\text{reçue}}} = 1 + \frac{C_v(T_1-T_4)}{C_p(T_3-T_2)} = 1 - \frac{T_4-T_1}{\gamma(T_3-T_2)} \]
 	En introduisant $x = V_1/V_2$ et $y = V_1/V_3$ :
 	\[ \eta = 1 - \frac{1}{\gamma x^{\gamma-1}} \left[ \frac{(x/y)^\gamma - 1}{x/y - 1} \right] \]
 	
 	\subsection*{3. Application numérique pour un moteur Diesel}
 	Caractéristiques : $x = 21$, $y = 7 \implies \frac{x}{y} = 3$. Avec $\gamma = 1,4$ :
 	\[ \eta = 1 - \frac{1}{1,4 \times (21)^{0,4}} \left( \frac{3^{1,4} - 1}{3 - 1} \right) \approx 1 - \frac{1}{1,4 \times 3,41} \left( \frac{4,65 - 1}{2} \right) \approx 0,617 \quad (61,7\%) \]
 	a) Rendement théorique : $61,7\%$.\\
 	b) Pour $N = 450\text{ tr/min}$ et consommation $c = 8\text{ L/100km}$ à $v = 147\text{ km/h}$.
 	Débit de carburant : $D = \frac{c \cdot v}{100} = 11,76\text{ L/h} \approx 3,27\cdot 10^{-3}\text{ L/s}$.
 	Masse injectée par cycle : $m_c = \frac{D \cdot \rho}{N} \approx 1,16\cdot 10^{-4}\text{ kg/cycle}$.\\
 	c) Puissance théorique maximale : $P_{\text{max}} = \eta \cdot D \cdot \rho \cdot q \approx 0,617 \times (3,27\cdot 10^{-6}\text{ m}^3\text{/s} \times 800\text{ kg/m}^3) \times 46,8\cdot 10^6\text{ J/kg} \approx 75,6\text{ kW}$.
 	
 	\newpage
 	
 	% ==========================================
 	\section*{Exercice 11 : }
 	% ==========================================
 	\subsection*{1. Représentation du cycle}
 	Tracé classique à deux isochores verticales ($2\rightarrow3$ et $4\rightarrow1$) et deux isentropiques courbes.
 	
 	\subsection*{2. Rendement théorique}
 	\[ \eta = 1 - \frac{T_4 - T_1}{T_3 - T_2} \]
 	Les transformations adiabatiques $1\rightarrow2$ et $3\rightarrow4$ donnent $\frac{T_1}{T_2} = \left(\frac{V_2}{V_1}\right)^{\gamma-1} = \frac{T_4}{T_3}$, d'où :
 	\[ \eta = 1 - \left(\frac{V_2}{V_1}\right)^{\gamma-1} = 1 - x^{1-\gamma} \]
 	
 	\subsection*{3. Application numérique}
 	Données : course $l = 10\text{ cm} = 0,1\text{ m}$, section $S = 50\text{ cm}^2 = 5\cdot 10^{-3}\text{ m}^2$, $V_2 = 100\text{ cm}^3 = 10^{-4}\text{ m}^3$.
 	Le volume balayé (cylindrée) est $V_b = S \cdot l = 5\cdot 10^{-4}\text{ m}^3 = 500\text{ cm}^3$.
 	Le volume maximal est $V_1 = V_2 + V_b = 600\text{ cm}^3$.
 	Le taux de compression volumétrique est donc $x = \frac{V_1}{V_2} = \frac{600}{100} = 6$.
 	\begin{itemize}
 		\item \textbf{a) Rendement théorique} ($\gamma = 1,4$) :
 		\[ \eta = 1 - 6^{-0,4} \approx 0,512 \quad (51,2\%) \]
 		\item \textbf{b) Température de fin de compression $T_2$ et travail} :
 		\[ T_2 = T_1 \cdot x^{\gamma-1} = 300 \times 6^{0,4} \approx 614,3\text{ K} \]
 		Avec $T_3 = 900\text{ K}$, la température de fin de détente $T_4$ est :
 		\[ T_4 = T_3 \cdot x^{1-\gamma} = 900 \times 6^{-0,4} \approx 439,5\text{ K} \]
 		Le travail reçu au cours d'un cycle par mole de fluide ($C_v = \frac{R}{\gamma-1}$) :
 		\[ W = \Delta U = C_v(T_2 - T_1) + C_v(T_4 - T_3) = \frac{R}{\gamma-1} (T_2 - T_1 + T_4 - T_3) \]
 		\[ W = \frac{8,314}{0,4} (614,3 - 300 + 439,5 - 900) \approx 20,78 \times (-146,2) \approx -3038\text{ J/mol} \]
 	\end{itemize}
 	
 	\subsection*{4. Comparaison avec un cycle de Carnot}
 	Pour $T_{\text{min}} = 300\text{ K}$ et $T_{\text{max}} = 900\text{ K}$, le rendement d'une machine de Carnot réversible est :
 	\[ \eta_C = 1 - \frac{T_{\text{min}}}{T_{\text{max}}} = 1 - \frac{300}{900} \approx 0,667 \quad (66,7\%) \]
 	Pour obtenir ce même rendement avec le cycle d'Otto :
 	\[ 1 - x^{1-\gamma} = 1 - \frac{T_1}{T_3} \implies x^{\gamma-1} = \frac{T_3}{T_1} = \frac{900}{300} = 3 \implies x = 3^{\frac{1}{0,4}} = 3^{2,5} \approx 15,6 \]
 	Le taux de compression volumétrique requis serait de $\approx 15,6$.
 	
 
	 
 	
	
	
	
	
	
	
	
	
	
	
	
\end{document}
